\documentclass[11pt]{article}
\usepackage{amsmath,amssymb,amsfonts,amsthm,graphicx,booktabs,hyperref}
\usepackage[margin=2.4cm]{geometry}
\usepackage{microtype}
\hypersetup{colorlinks=true,linkcolor=blue,citecolor=blue,urlcolor=blue}
\newtheorem{theorem}{Theorem}
\newtheorem{lemma}{Lemma}
\newtheorem{proposition}{Proposition}
\newtheorem{corollary}{Corollary}
\theoremstyle{definition}
\newtheorem{definition}{Definition}
\title{Finite-Donor Primitive Fluxes\\ and the Nonexistence of a Closed Rest Point}
\author{Your Name}
\date{\today}
\begin{document}
\maketitle

\begin{abstract}
The working equilibrium of the companion delay-differential core is a frozen-donor quasi-equilibrium: it is sustained by a one-way geological transfer of order $4.652$ stock units per year and is not a rest point of any closed finite-donor system. This paper replaces the derived target that manufactured that quasi-equilibrium with a primitive-flux ledger in which every transfer is donor-limited and every compartment, including $A^{\mathrm{geo}}$, is a state. The natural block $N+A^{\mathrm{act}}+A^{\mathrm{geo}}+U$ loses mass exactly at the extraction and mining rates. The only rest points with vanishing extraction are extinction plus geochemical equilibrium. There is no interior rest point at positive effort. We do not claim finite-time tracking of the companion working core by the primitive system. The primitive field differs from the working field by an $O(1)$ term at the working point; the companion's $\varepsilon_G$ estimate belongs to a different derived-target completion.
\end{abstract}

\section{Relation to the companion}
\label{sec:relation}

The companion manuscript~\cite{paperI} tracks a mass-conserved vector ledger and reduces, under institutional failure, to an effort-bounded delay-differential core. Its working four-state equilibrium $(N^*,A^{\mathrm{act}*},Z^*,E^*)=(89.526,397.87,\ln 2/10,2.090)$ is held by the \emph{derived} target
\[
A^{\mathrm{eq}}=A^{\mathrm{eq,intrinsic}}+\frac{\kappa_A K-\gamma_U U}{\omega_A},
\]
which cancels detritus return and produces a high-$A$ rest of the active pool while $A^{\mathrm{geo}}$ is held fixed. The geological support required is
\begin{equation}
\omega_A\bigl(A^{\mathrm{eq,W}}-A^{\mathrm{act}*}\bigr)=4.652133\ldots
\label{eq:flux}
\end{equation}
stock units per year. That number is not a rest-point residual of a closed $A^{\mathrm{geo}}$ equation. It is the flux that must be supplied, every year, by a donor the companion treats as a parameter.

Here the donor is a state and every transfer is a primitive. The delays $\tau_\pm$ of~\cite{paperI} remain those of the frozen-donor specialisation.

\section{Primitive ledger}
\label{sec:ledger}

Let $s=A^{\mathrm{act}}/(A^{\mathrm{act}}+A_0)$ and $\sigma=A^{\mathrm{geo}}/(A^{\mathrm{geo}}+A_{g0})$. Net regeneration and gross uptake are the companion's constitutive laws
\begin{equation}
R(N,A^{\mathrm{act}})=rN\Bigl(1-\frac{N}{K}\Bigr)s,\qquad
B=R+\kappa_A N s,\qquad
T=\kappa_A N s.
\label{eq:RB}
\end{equation}
The four primitives involving the donor are
\begin{equation}
\begin{aligned}
e_{GA}&=\omega_A\bigl[A^{\mathrm{eq,intrinsic}}\bigr]_+\sigma,\\
e_{AG}&=\omega_A A^{\mathrm{act}},\\
C^{A,\mathrm{lim}}&=C^A\sigma,\\
\gamma_U U&=\text{detritus return}.
\end{aligned}
\label{eq:prim}
\end{equation}
No derived target appears. Recharge cannot run backward: $e_{GA}=0$ at $A^{\mathrm{geo}}=0$ and at a nonpositive intrinsic target. Mining is donor-limited in the same way extraction is.

The closed natural block, under the companion's institutional-failure specialisation $\mu=\nu=\rho=0$ and $C^A=0$, is
\begin{equation}
\begin{aligned}
\dot N&=R(N,A^{\mathrm{act}})-qEN,\\
\dot A^{\mathrm{act}}&=-B+e_{GA}-e_{AG}+\gamma_U U,\\
\dot A^{\mathrm{geo}}&=-e_{GA}+e_{AG},\\
\dot U&=T-\gamma_U U,
\end{aligned}
\label{eq:closed}
\end{equation}
together with the companion's memory--effort pair $(Z,E)$ driven by $qEN-R$ (never by mining). Parameters are those of Candidate~A in~\cite{paperI}: $r=0.02$, $K=100$, $q=0.001$, $\kappa_A=0.05$, $\omega_A=10^{-3}$, $A_0=1$, $A^{\mathrm{eq,intrinsic}}=50$, $\gamma_U=0.2$.

\section{Conservation and invariance}

\begin{theorem}[Natural-block mass identity]
\label{thm:mass}
Along every absolutely continuous solution of~\eqref{eq:closed},
\[
\frac{d}{dt}\bigl(N+A^{\mathrm{act}}+A^{\mathrm{geo}}+U\bigr)=-qEN.
\]
If product, waste, and the inert sink are restored with the companion's donor-limited routing, the full ledger $N+A^{\mathrm{act}}+A^{\mathrm{geo}}+U+P+W+I$ is constant.
\end{theorem}

\begin{proof}
Adding the four equations in~\eqref{eq:closed} gives
\begin{align*}
\dot N+\dot A^{\mathrm{act}}+\dot A^{\mathrm{geo}}+\dot U
&=(R-qEN)+(-B+e_{GA}-e_{AG}+\gamma_U U)+(-e_{GA}+e_{AG})+(T-\gamma_U U).
\end{align*}
By~\eqref{eq:RB}, $B=R+T$, so $R-B+T=0$. The pairs $\pm e_{GA}$, $\pm e_{AG}$, and $\pm\gamma_U U$ cancel, and the remainder is $-qEN$. If product, waste, and the inert sink are restored with the donor-limited routing of~\cite{paperI}, each extracted unit appears in $P$ or $W$ and each recovery from $W$ returns to $A^{\mathrm{act}}$, so the seven-compartment sum is constant.
\end{proof}

\begin{theorem}[Orthant invariance]
\label{thm:orthant}
The nonnegative orthant in $(N,A^{\mathrm{act}},A^{\mathrm{geo}},U,Z,E)$ is forward invariant. At $A^{\mathrm{geo}}=0$ one has $e_{GA}=0$ and $\dot A^{\mathrm{geo}}=e_{AG}\ge 0$. At $A^{\mathrm{act}}=0$ one has $R=B=T=e_{AG}=0$ and $\dot A^{\mathrm{act}}=e_{GA}+\gamma_U U\ge 0$. At $N=0$ extraction and uptake vanish.
\end{theorem}

\begin{proof}
The right-hand side of~\eqref{eq:closed} is locally Lipschitz on a neighbourhood of the closed orthant (each Michaelis--Menten factor $s$, $\sigma$ is $C^\infty$ for nonnegative arguments after the usual $C^1$ extension through the origin). On the face $A^{\mathrm{geo}}=0$ one has $\sigma=0$, hence $e_{GA}=0$ and $\dot A^{\mathrm{geo}}=e_{AG}=\omega_A A^{\mathrm{act}}\ge 0$. On $A^{\mathrm{act}}=0$ one has $s=0$, so $R=B=T=e_{AG}=0$ and $\dot A^{\mathrm{act}}=e_{GA}+\gamma_U U\ge 0$. On $N=0$ one has $R=T=qEN=0$ and $\dot N=0$. On $U=0$, $\dot U=T\ge 0$. Nagumo's inward-pointing criterion \cite{nagumo1942} therefore yields forward invariance of the orthant.
\end{proof}

\section{Rest points}

\begin{theorem}[No interior rest at positive effort]
\label{thm:no-rest}
Suppose $E\equiv E_*>0$ is constant. A rest point of~\eqref{eq:closed} satisfies $R+C^{A,\mathrm{lim}}=0$ after restoring optional mining. With $C^A=0$ this is $R=0$, hence $N=0$ or $N=K$ or $A^{\mathrm{act}}=0$. None of these is compatible with $E_*>0$ and $N_*>0$:
\begin{enumerate}
\item $N=K$ and $E_*>0$ give $\dot N=-qE_*K<0$.
\item $A^{\mathrm{act}}=0$ and $A^{\mathrm{geo}}>0$ give $\dot A^{\mathrm{act}}=\omega_A A^{\mathrm{eq,intrinsic}}\sigma>0$.
\item $N=0$ forces $R=T=0$ and reduces to the extinction--geochemical rest of Theorem~\ref{thm:extinction}.
\end{enumerate}
In particular the companion's working point $(N^*,A^{\mathrm{act}*})=(89.526,397.87)$ is not a rest point of~\eqref{eq:closed} at $E=E^*\approx 2.090$.
\end{theorem}

\begin{proof}
At a rest point, $\dot U=0$ forces $\gamma_U U=T$. Adding $\dot A^{\mathrm{act}}+\dot A^{\mathrm{geo}}$ then gives $-B+\gamma_U U=0$, hence $B=T$. But $B=R+T$, so $R=0$. From~\eqref{eq:RB}, $R=rN(1-N/K)s=0$ implies $N=0$ or $N=K$ or $s=0$ (i.e.\ $A^{\mathrm{act}}=0$).

If $N=K$ and $E_*>0$, then $\dot N=0-qE_*K<0$, contradicting rest. If $A^{\mathrm{act}}=0$ and $A^{\mathrm{geo}}>0$, then $\sigma>0$ and $\dot A^{\mathrm{act}}=\omega_A A^{\mathrm{eq,intrinsic}}\sigma>0$, again not rest. If $N=0$, then $R=T=0$ and the remaining equations reduce to Theorem~\ref{thm:extinction}.

At the working point of~\cite{paperI} one has $R^*=qE^*N^*\approx 0.187>0$, which contradicts $R=0$.
\end{proof}

\begin{theorem}[Extinction--geochemical rest]
\label{thm:extinction}
The set $N=0$, $U=0$, $E$ arbitrary, $A^{\mathrm{act}}=A^{\mathrm{eq,intrinsic}}\sigma$, $\sigma=A^{\mathrm{geo}}/(A^{\mathrm{geo}}+A_{g0})$ consists of rest points of the natural block. If $A_{g0}=0$ and $\sigma\equiv 1$ for $A^{\mathrm{geo}}>0$, this is the ray $A^{\mathrm{act}}=A^{\mathrm{eq,intrinsic}}$, $A^{\mathrm{geo}}\ge 0$.
\end{theorem}

\begin{proof}
Set $N=U=0$. Then $R=B=T=0$, $e_{GA}=\omega_A A^{\mathrm{eq,intrinsic}}\sigma$, $e_{AG}=\omega_A A^{\mathrm{act}}$, and
\begin{align*}
\dot A^{\mathrm{act}}&=\omega_A A^{\mathrm{eq,intrinsic}}\sigma-\omega_A A^{\mathrm{act}},\\
\dot A^{\mathrm{geo}}&=-\omega_A A^{\mathrm{eq,intrinsic}}\sigma+\omega_A A^{\mathrm{act}}.
\end{align*}
Both vanish if and only if $A^{\mathrm{act}}=A^{\mathrm{eq,intrinsic}}\sigma$. If $A_{g0}=0$ and $A^{\mathrm{geo}}>0$, then $\sigma=1$ and the rest set is the ray $A^{\mathrm{act}}=A^{\mathrm{eq,intrinsic}}$, $A^{\mathrm{geo}}\ge 0$.
\end{proof}

Institutional memory of~\cite{paperI} yields $E\to E^*$ at $N=0$. Extraction then vanishes identically, so the point is consistent with Theorem~\ref{thm:extinction} and is not an interior rest.

\section{Relation to the working core on finite horizons}

The companion's frozen-donor working core and the primitive system~\eqref{eq:closed} are different completions. The working completion uses the derived target $A^{\mathrm{eq,W}}$, whereas~\eqref{eq:closed} recharges toward $A^{\mathrm{eq,intrinsic}}$. At the companion working point, the primitive and working $A^{\mathrm{act}}$ fields therefore differ by an $O(1)$ term even when $\sigma\approx1$.

Consequently, the quantity
\[
\varepsilon_G(T)=G_0^{-1}\int_0^T|e_{GA}-e_{AG}|\,dt
\]
can quantify cumulative donor draw in the derived-target completion, but it is not a tracking error bound between~\eqref{eq:closed} and the working core. A finite-time comparison for the primitive system requires a separate continuous-dependence estimate based on its actual vector field and initial data; no such estimate is asserted here.

The working flux value $4.652133$ stock units per year remains useful as a diagnostic: it shows directly that the frozen-donor working point requires sustained geological support. It is not a primitive-system equilibrium flux.

\section{Long-time donor decline}

\begin{theorem}[Integrable extraction]
\label{thm:integrable}
Let $M=N+A^{\mathrm{act}}+A^{\mathrm{geo}}+U$. Then $M(t)=M(0)-\int_0^t qE(s)N(s)\,ds\ge 0$, so
\[
\int_0^\infty qE(s)N(s)\,ds\le M(0)<\infty.
\]
In particular $qEN\in L^1(0,\infty)$. There is no trajectory of~\eqref{eq:closed} on which extraction remains at the companion's working value $qE^*N^*\approx 0.187$ for all time.
\end{theorem}

\begin{proof}
By Theorem~\ref{thm:mass}, $M(t)=M(0)-\int_0^t qE(s)N(s)\,ds$. Forward invariance (Theorem~\ref{thm:orthant}) gives $M(t)\ge 0$, hence the improper integral of $qEN$ is at most $M(0)$. If $qEN\equiv qE^*N^*\approx 0.187$ for all $t\ge 0$, the integral would diverge, a contradiction.
\end{proof}

The companion's $\tau_+\approx 150$~yr upper cycle is therefore a frozen-donor object. On the closed system it can persist only on a finite donor-budget timescale; under a sustained lower extraction flux $c>0$, the duration is bounded above by $G_0/c$, and an $O(G_0/\text{flux})$ scale is the appropriate heuristic. At $G_0/A^{\mathrm{act}*}=10^3$ that length is tens of thousands of years, far above the institutional delays of~\cite{paperI}. Whether the frozen-donor local Hopf structure persists as a slowly drifting transient in the closed donor system is an open slow-passage/adiabatic-continuation problem; the mass budget alone does not prove it.

\section{Frozen-donor limit and its scope}

Rescaling $G=G_0g$ with $g(0)=1$ gives
\[
\dot g=-G_0^{-1}(e_{GA}-e_{AG}).
\]
For the primitive system~\eqref{eq:closed}, the limit $G_0\to\infty$ freezes $g$ but does not restore the companion's derived target: the limiting recharge field still uses $A^{\mathrm{eq,intrinsic}}$, not $A^{\mathrm{eq,W}}$. Consequently this scaling is not a regular perturbation of the working four-state vector field. Local Hopf persistence of the working core under this primitive scaling is not claimed. A different derived-target completion would be required before a regular-perturbation theorem could be formulated.

\section{Discussion}

With $G(t)$ included as a state,~\eqref{eq:closed} is an autonomous RFDE with a slow donor coordinate. Periodic orbits of the frozen-donor core can persist only as transients on the finite donor budget; the precise transient duration and any continuation of attractors in $(G_0,\tau)$ require a separate computation.

\section{Conclusion}

The working point of~\cite{paperI} is not a rest point of the closed primitive-flux system. The primitive system has a finite donor budget and extraction is integrable, so a constant working flux cannot be maintained indefinitely. No $O(\varepsilon_G(T))$ tracking theorem for the primitive completion is claimed.

\begin{thebibliography}{9}
\bibitem{paperI}
Y.\ Name, ``Scarcity-driven capital liquidation and delay-amplified instability: a vector-valued flow-balance framework for sustainability,'' companion manuscript, 2026.
\bibitem{hale1993}
J.\ K.\ Hale and S.\ M.\ Verduyn Lunel, \emph{Introduction to Functional Differential Equations}, Springer, 1993.
\bibitem{nagumo1942}
M.\ Nagumo, \"Uber die Lage der Integralkurven gew\"ohnlicher Differentialgleichungen, \emph{Proc.\ Phys.-Math.\ Soc.\ Japan} 24 (1942), 551--559.
\end{thebibliography}

\end{document}
